% ===============================================================
% 系统实现
% ===============================================================
\section{系统实现}

% ---------------------------------------------------------------
\begin{frame}{\code{/dev/kcov} 和 \code{/proc/kallsyms}}
\begin{columns}[T,onlytextwidth]
\begin{column}{0.64\textwidth}
\footnotesize
\renewcommand{\arraystretch}{1.3}
\begin{tabularx}{\linewidth}{@{}>{\ttfamily}l>{\raggedright\arraybackslash}X@{}}
  \toprule
  \normalfont\textbf{接口} & \textbf{语义与防御规则} \\
  \midrule
  open & 只允许 root（否则 \code{EPERM}），每次打开分配独立的 \code{KcovFile} \\
  KCOV\_INIT\_TRACE & 条目数 $\in[2,2^{20}]$，否则 \code{EINVAL}；重复初始化返回 \code{EBUSY} \\
  KCOV\_ENABLE & 参数必须为 0（仅支持 trace-pc）；未初始化 \code{EINVAL}，已绑定 \code{EBUSY} \\
  KCOV\_DISABLE & 只有绑定该缓冲区的任务可以解绑，否则 \code{EINVAL} \\
  mmap & 必须 \code{MAP\_SHARED}、偏移为 0、长度恰好等于缓冲区大小 \\
  read/write/lseek & 不支持流式读写与寻道，返回 \code{EINVAL} / \code{ESPIPE} \\
  /proc/kallsyms & 只允许 root 读取，导出链接脚本锚定的 \code{\_stext}、\code{\_etext} \\
  \bottomrule
\end{tabularx}
\end{column}
\begin{column}{0.33\textwidth}
\begin{tcolorbox}[rustcard,title={回调路径零分配},fontupper=\footnotesize]
  初始化时一次性分配并\textbf{预先提交}全部物理页，记录 PC 时不缺页、不分配内存
\end{tcolorbox}
\vspace{1mm}
\begin{tcolorbox}[tcbcard,title={符号可追溯},fontupper=\footnotesize]
  镜像脚本保留未剥离的 \code{vmlinux}，并校验其中有 \code{.symtab}、
  \code{.debug\_info} 和覆盖回调符号
\end{tcolorbox}
\vspace{1mm}
\begin{tcolorbox}[syzcard,title={明确拒绝},fontupper=\footnotesize]
  比较操作数追踪、远程覆盖等扩展模式
\end{tcolorbox}
\end{column}
\end{columns}
\end{frame}

% ---------------------------------------------------------------
\begin{frame}{\syz{} 适配}
\begin{tcbraster}[raster columns=3,raster equal height,raster column skip=2mm]
\begin{tcolorbox}[syzcard,title={执行器}]
  \code{debugfs} 不存在时回退到 \code{/dev/kcov}；放宽可选特性检查
\end{tcolorbox}
\begin{tcolorbox}[syzcard,title={符号解析}]
  回退到 \code{\_\_stext}；无覆盖对照组可缺失 \code{kallsyms}
\end{tcolorbox}
\begin{tcolorbox}[syzcard,title={虚拟机管理}]
  \aster{} 的 QEMU 配置；快照与自适应生命周期管理
\end{tcolorbox}
\end{tcbraster}

\vspace{1mm}
\begin{tcolorbox}[syzcard,title={Rust panic 报告解析}]
\begin{minipage}[c]{0.65\linewidth}
  \begin{tcolorbox}[terminal]
\textcolor{SATcb!60}{[\ \ 91.693]} ERROR: Uncaught panic:\\
\hspace*{1.5em}not implemented\\
\hspace*{1.5em}at /src/asterinas/kernel/core/src/syscall/timer\_create.rs:74\\
Printing stack trace: ...
  \end{tcolorbox}
\end{minipage}\hfill
{\large\textcolor{SASyz}{$\Rightarrow$}}\hfill
\begin{minipage}[c]{0.3\linewidth}
  \begin{tcolorbox}[terminal,colback=SASyz,colframe=SASyz,left=3pt,right=3pt]
asterinas panic: not implemented at src/syscall/\allowbreak timer\_create.rs:LINE
  \end{tcolorbox}
\end{minipage}

\vspace{1mm}
路径只保留最后 3 级、行号归一化 $\Rightarrow$ 不同构建路径和代码行漂移下\textbf{仍能正确去重}
\end{tcolorbox}

\vspace{1mm}
{\scriptsize\muted{\faIcon{wrench}\ 移植中另修复了 \aster{} 的问题（\textbf{非}模糊测试发现）：
NixOS 镜像安装与串口控制台、\code{PTRACE\_TRACEME} 在 Yama 下重复加锁父进程导致的死锁、
\code{sys\_set\_tid\_address} 中的 \code{todo!()} panic，均附带回归测试}}
\end{frame}

% ---------------------------------------------------------------
\begin{frame}{Rust 工具链和 TUI}
\begin{columns}[T,onlytextwidth]
\begin{column}{0.5\textwidth}
\small
\textbf{统一入口 \code{bin/syzaster}}（只依赖 \code{std}）

\vspace{1mm}
\renewcommand{\arraystretch}{1.2}
\begin{tabularx}{\linewidth}{@{}>{\ttfamily\color{SARust}}l>{\raggedright\arraybackslash}X@{}}
  check-inst & 插桩边界校验 \\
  check-kcov & 在 VM 中运行 kcov ABI 测试 \\
  kcov-rustc & \code{KCOV=1} 构建用的编译器包装器 \\
  eval-fuzz & 统计评估：$\hat{A}_{12}$、Mann-Whitney U \\
  tui & 键盘驱动的终端界面 \\
\end{tabularx}

\vspace{1mm}
\textbf{Shell 脚本}：打补丁、镜像、QEMU、分诊复现

\textbf{测试配置}：全内核、4 个子系统、无覆盖对照
\end{column}
\begin{column}{0.47\textwidth}
\centering
\imageorplaceholder{images/SyzAster.png}{\linewidth}{4.6cm}{\code{syzaster tui} 监控页}

\vspace{1mm}
{\scriptsize 内核配置 · 模糊测试配置 · 监控（覆盖、崩溃、执行速度曲线与问题列表）}
\end{column}
\end{columns}

\vspace{2mm}
\begin{tcbraster}[raster columns=5,raster equal height,raster column skip=2mm]
\stattile{rustcard}{0}{外部 crate 依赖}
\stattile{rustcard}{47}{Rust 测试通过}
\stattile{syzcard}{60}{Go 单元测试}
\stattile{tcbcard}{2}{\aster{} 版本}
\stattile{accentcard}{\faDocker}{一站式容器}
\end{tcbraster}
\end{frame}

% ---------------------------------------------------------------
\begin{frame}{Rust 特性}
\small
\renewcommand{\arraystretch}{1.45}
\begin{tabularx}{\textwidth}{@{}>{\raggedright}p{3.7cm}>{\raggedright}p{4.4cm}
    >{\raggedright\arraybackslash}X@{}}
  \toprule
  \textbf{语言特性} & \textbf{在 \sa{} 中的位置} & \textbf{带来的保证} \\
  \midrule
  \rust{裸函数}\newline\code{\#[unsafe(naked)]}、\code{naked\_asm!}
    & 覆盖回调跳板 \code{\_\_sanitizer\_cov\_trace\_pc}
    & 汇编只有 2 条指令，满足 SysV 栈对齐，其余逻辑回到安全 Rust \\
  \rust{\code{Option} + \code{try\_lock()}}
    & \code{record\_pc} 的非重入记录
    & 编译器强制处理“拿不到锁”的分支，不会自旋死锁 \\
  \rust{RAII / \code{Drop}}
    & \code{Attachment} 自动解绑；本地与远程临时目录自动清理
    & 任何退出路径（含异常退出）都能确定地释放资源 \\
  \rust{\code{Result} + \code{read\_once}/\code{write\_once}}
    & \code{KcovBuffer::record}
    & 没有裸指针运算，每次共享内存访问的错误都必须处理 \\
  \rust{切片边界检查}
    & 自研 ELF 解析器 \code{src/pkg/elf.rs}
    & 畸形 ELF 只会得到解析错误，不会越界读取 \\
  \bottomrule
\end{tabularx}
\end{frame}
